% GENERATED FILE. Edit canonical lessons, then run npm run generate:books.
% canonical-source-hash: fnv1a64:e607a0cb91dc72e3
% canonical-lessons: PA-C206-go-return, PA-C206-gae, PA-C206-e, PA-C206-gae-review, PA-C206-they-went, PA-C206-payoff

\chapter{They Went}
\label{ch:pa-a2-gae}
\hlchaptermodality{206}

\begin{chapteropening}
\textbf{By the end of this chapter:} \emph{I can recognise and use one bounded masculine plural form for went.}

Recognise and produce one bounded they-went frame through listening, speaking, reading, and writing.
\end{chapteropening}

\section[\texorpdfstring{\pa{ਉਹ} · \pa{ਜਾਣਾ} (oh; jaana)}{ਉਹ · ਜਾਣਾ (oh; jaana)}]{Return to a familiar pronoun and verb}

\label{lesson:PA-C206-go-return}



\begin{quote}
\emph{Hear:} \emph{oh; jaana}; YOU SAY: \textbf{\pa{ਉਹ}; \pa{ਜਾਣਾ}}
\end{quote}

\begin{grammarlens}[title={Grammar Lens}]
Keep the familiar pronoun and the familiar infinitive ready. The next lesson will own one new past form; this opener does not reveal it.
\end{grammarlens}

\subsection*{Your turn}
\begin{itemize}
\raggedright
  \item \emph{Hear:} \emph{they or respectful he}; YOU SAY: \textbf{\pa{ਉਹ}}
  \item \emph{Hear:} \emph{to go}; YOU SAY: \textbf{\pa{ਜਾਣਾ}}
  \item Say the familiar pronoun, then the familiar infinitive.
\end{itemize}

\begin{tcolorbox}[breakable,colback=teal!4,colframe=teal!35!black,title={Before you move on}]
Produce the two familiar forms from meaning. This opener owns no new form.
\end{tcolorbox}
\section[\texorpdfstring{\pa{ਗਏ} (gae)}{ਗਏ (gae)}]{One bounded past form of to go}

\label{lesson:PA-C206-gae}



\begin{quote}
Retrieve the familiar infinitive \textbf{\pa{ਜਾਣਾ}}, “to go,” then listen for one irregular past form.
\end{quote}

\begin{grammarlens}[title={Grammar Lens: direct owner}]
\begin{quote}
\textbf{\pa{ਗਏ}} — \emph{gae} — went, masculine plural
\end{quote}

Hear and say the form before reading it. Punjabi also uses this masculine plural form for one respectfully mentioned masculine person. Keep the irregular form whole. This chapter practises only the small “they went” frame; it does not open the wider perfective or agreement system.
\end{grammarlens}

\begin{cousinweb}
Punjabi University Patiala's lexical statistics identifies \textbf{\pa{ਗਏ}} as the masculine plural perfect form of the verb “go.” Its verb reference places the same form in the perfective family of that verb, and its reference grammar independently labels it masculine plural. Sources: \href{https://www.learnpunjabi.org/statistics.html}{Punjabi University lexical statistics}, \href{https://www.learnpunjabi.org/Verb.html}{Punjabi University verb reference}, and \href{https://pt.learnpunjabi.org/assets/a\%20reference\%20grammar\_final.pdf}{Punjabi University reference grammar}.

Those sources support the modern form. Reuse the history already attached to \textbf{\pa{ਜਾਣਾ}}; do not invent a second etymology for an inflected form. For this lesson, remember the whole irregular form and one bounded meaning.
\end{cousinweb}

\subsection*{Your turn}
\begin{itemize}
\raggedright
  \item \emph{Hear:} \emph{gae}; YOU SAY: \textbf{\pa{ਗਏ}}
  \item \emph{Read it:} \textbf{\pa{ਗਏ}}
  \item \emph{Write it:} copy \textbf{\pa{ਗਏ}} once, checking the whole form
\end{itemize}

\begin{tcolorbox}[breakable,colback=teal!4,colframe=teal!35!black,title={Before you move on}]
Produce the masculine plural form for “went.”
\end{tcolorbox}
\section[\texorpdfstring{\pa{ਏ} (e)}{ਏ (e)}]{\pa{ਏ} — one vowel from the went form you can say}

\label{lesson:PA-C206-e}



\begin{quote}
Say \textbf{\pa{ਗਏ}}, the whole “went” form from the previous lesson.
\end{quote}

\subsection*{Script you'll notice: \pa{ਏ}}
\begin{quote}
\textbf{\pa{ਗ}} + \textbf{\pa{ਏ}} $\to$ \textbf{\pa{ਗਏ}}
\end{quote}

\begin{quote}
\textbf{\pa{ਏ}} — \emph{e} — an independent Gurmukhi vowel
\end{quote}

Keep \textbf{\pa{ਏ}} in view. It is the second letter in the known word \textbf{\pa{ਗਏ}}.

No stroke order is taught here: the track does not yet have a verified source for one.

\subsection*{Writing: copy what you can see}
Keep \textbf{\pa{ਏ}} in front of you. Trace its printed shape once with a finger, then copy it once with a pencil while the model stays visible. Say \emph{e} when you finish.

\begin{quote}
This is observe-and-copy with the model visible, not writing from memory.
\end{quote}

\subsection*{Your turn}
\begin{itemize}
\raggedright
  \item \emph{Look:} at \textbf{\pa{ਏ}}; YOU SAY: \emph{e}
  \item \emph{Read it:} \textbf{\pa{ਗਏ}}, then point to \textbf{\pa{ਏ}}
  \item \emph{Write it:} copy \textbf{\pa{ਏ}} once while the model remains visible
\end{itemize}

\begin{tcolorbox}[breakable,colback=teal!4,colframe=teal!35!black,title={Before you move on}]
Point to \textbf{\pa{ਏ}}, say \emph{e}, then say the whole word \textbf{\pa{ਗਏ}}.

Source: \href{https://www.unicode.org/charts/PDF/U0A00.pdf}{Unicode Gurmukhi chart}.
\end{tcolorbox}
\section[\texorpdfstring{\pa{ਗਏ} (gae)}{ਗਏ (gae)}]{Retrieve the new went form on its own}

\label{lesson:PA-C206-gae-review}



\begin{quote}
\emph{Hear:} \emph{gae}; YOU SAY: \textbf{\pa{ਗਏ}}, then point to \textbf{\pa{ਏ}}
\end{quote}

\subsection*{Your turn}
\begin{itemize}
\raggedright
  \item \emph{Hear:} \emph{went, masculine plural}; YOU SAY: \textbf{\pa{ਗਏ}}
  \item \emph{Read it:} \textbf{\pa{ਗਏ}}, then point to \textbf{\pa{ਏ}} and say \emph{e}
  \item \emph{Write it:} cover and write \textbf{\pa{ਗਏ}}
\end{itemize}

\begin{grammarlens}[title={Grammar Lens}]
Keep one whole irregular form and one bounded past meaning here. The familiar pronoun returns only after this separate retrieval.
\end{grammarlens}

\begin{tcolorbox}[breakable,colback=teal!4,colframe=teal!35!black,title={Before you move on}]
Say and then write the form once from memory.
\end{tcolorbox}
\section[\texorpdfstring{\pa{ਉਹ} \pa{ਗਏ} (oh gae)}{ਉਹ ਗਏ (oh gae)}]{Put the new went form after a familiar pronoun}

\label{lesson:PA-C206-they-went}



\begin{quote}
Retrieve the familiar pronoun for “they,” then the new masculine plural form for “went.”
\end{quote}

\begin{grammarlens}[title={Grammar Lens}]
\begin{quote}
\textbf{\pa{ਉਹ} \pa{ਗਏ}} — \emph{oh gae} — they went
\end{quote}

For this one frame, put familiar \textbf{\pa{ਉਹ}} before the new whole form \textbf{\pa{ਗਏ}}. In a different context the same masculine plural form can respectfully refer to one masculine person. Nothing else in the perfective or agreement system is added.
\end{grammarlens}

\subsection*{Your turn}
\begin{itemize}
\raggedright
  \item \emph{Hear:} \emph{oh gae}; YOU SAY: \textbf{\pa{ਉਹ} \pa{ਗਏ}}
  \item \emph{Read it:} \textbf{\pa{ਉਹ} \pa{ਗਏ}}
  \item \emph{Write it:} cover and write \textbf{\pa{ਉਹ} \pa{ਗਏ}}
\end{itemize}

\begin{tcolorbox}[breakable,colback=teal!4,colframe=teal!35!black,title={Before you move on}]
Produce the complete “they went” frame.
\end{tcolorbox}
\section[\texorpdfstring{\pa{ਉਹ} \pa{ਗਏ} (oh gae)}{ਉਹ ਗਏ (oh gae)}]{A bounded they-went payoff}

\label{lesson:PA-C206-payoff}



\begin{quote}
Retrieve the familiar infinitive “to go,” then the new whole form for masculine plural “went.”
\end{quote}

\subsection*{Your turn: four skills}
\begin{itemize}
\raggedright
  \item \emph{Hear:} \emph{oh gae}; YOU SAY: \textbf{\pa{ਉਹ} \pa{ਗਏ}}
  \item \emph{Read it:} \textbf{\pa{ਉਹ} \pa{ਗਏ}}
  \item \emph{Write it:} cover the line and write \textbf{\pa{ਉਹ} \pa{ਗਏ}}
\end{itemize}

\begin{tcolorbox}[breakable,colback=teal!4,colframe=teal!35!black,title={Before you move on}]
Say the new form, then produce the complete familiar-pronoun frame.

This is the twenty-ninth prerequisite tranche exposed by Punjabi A2 mock drafting. It does not unblock Punjabi A2 mock 1 by itself. It is not an A2 mock, not A2 readiness or attainment, and not external or human validation.
\end{tcolorbox}
